% Generated draft: requires mathematical and source review.
\documentclass[11pt]{article}
\usepackage[margin=25mm]{geometry}
\usepackage{fontspec}
\setmainfont{Latin Modern Roman}
\usepackage{amsmath,amssymb,mathtools,unicode-math}
\setmathfont{Latin Modern Math}
\usepackage{enumitem,needspace,xurl,hyperref}
\setlength{\parindent}{0pt}
\setlength{\parskip}{5pt}
\setlength{\emergencystretch}{3em}
\newcommand{\N}{\mathbb{N}}
\newcommand{\Z}{\mathbb{Z}}
\newcommand{\Q}{\mathbb{Q}}
\newcommand{\R}{\mathbb{R}}
\newcommand{\C}{\mathbb{C}}
\newcommand{\F}{\mathbb{F}}
\newcommand{\A}{\mathbb{A}}
\newcommand{\PP}{\mathbb{P}}
\newcommand{\E}{\mathbb{E}}
\newcommand{\eps}{\varepsilon}
\newcommand{\dd}{\,\mathrm{d}}
\begin{document}
\section*{UnsolvedMath ID: 701580}
\textbf{Source problem code:} catalogue entry
\subsection{Definitions and mathematical statement}
Let \(p\geq 673\) be a prime and let \(r\geq 2\) be an integer. Let \(F_r\) be the free group with free generators \(a_1,\ldots,a_r\), and put
\[
B=B(r,p):=F_r/\langle\!\langle g^p:g\in F_r\rangle\!\rangle,
\]
where the denominator denotes normal closure. Write \(a_1,\ldots,a_r\) also for their images in \(B\). Thus \(B\) is the free group of rank \(r\) in the variety of groups satisfying \(y^p=1\).

A one-variable equation over \(B\) is an equation \(w(a_1,\ldots,a_r,x)=1\), where \(w\) is a group word with coefficients in \(B\), represented by an element of the free product \(B*\langle x\rangle\), with \(\langle x\rangle\cong\Z\). For \(b\in B\), evaluation fixes the coefficients and substitutes \(b\) for \(x\). Define
\[
S_w:=\{b\in B:w(a_1,\ldots,a_r,b)=1\}.
\]
Interpreting the source's term ``non-identical'' as meaning that the equation does not hold for every substitution in \(B\), the question is whether, for every such \(p,r\) and every such word \(w\),
\[
\bigl(\exists b\in B\; w(a_1,\ldots,a_r,b)\neq 1\bigr)
\quad\Longrightarrow\quad |S_w|<\infty.
\]
The source does not explicitly define ``non-identical''. The displayed formulation uses the interpretation just stated; whether the intended condition instead refers to nontriviality of a word in a relatively free Burnside group requires confirmation. Formal nontriviality in \(B*\langle x\rangle\) alone is a different condition. No uniform bound on the finite number of solutions is requested.
\subsection{Short English statement}
For every prime \(p\geq 673\) and integer \(r\geq 2\), must every one-variable equation over the free Burnside group \(B(r,p)\) that is not satisfied by all elements have only finitely many solutions? This formulation interprets the source's undefined term ``non-identical'' as failure under at least one substitution; the intended meaning requires review.
\subsection{Sources}
\begin{itemize}
\item Supplied local UnsolvedMath registry record, permanent ID \(701580\), exact source problem code \texttt{catalogue entry}, title ``Kourovka Notebook Problem 19.32''; upstream revision \texttt{d371ee917ac9647fbb06cd6fc529deec6f3dda82}.
\item The supplied record attributes Problem 19.32 to A. Gaglione, proposed in 2018, and identifies its source as \emph{The Kourovka Notebook (Unsolved Problems in Group Theory)}, 19th issue (2018), with the 21st issue edition (2026), \texttt{arXiv:1401.0300v48}. The supplied collection description names E. I. Khukhro and V. D. Mazurov as editors.
\item The registry labels the problem ``open'' and records an update dated 2026-10-16. These are supplied metadata, not an independently verified assessment of the literature or mathematical status. The cited source has not been independently consulted for this draft.
\end{itemize}
\end{document}
